% Copyright (c) 2026 Marcio Luiz Lima de Oliveira and V. el Hakim.
% Distribution and modification are subject to LPPL version 1.3c.
% SPDX-License-Identifier: LPPL-1.3c
% Part of the LPPL-covered work listed in MANIFEST.txt.
% See LICENSE and LICENSING.txt for the licence and scope.

\chapter{Start with the structure}\label{chap:01}
\renewcommand{\FigLocation}{Chapters/Chapter_1/fig/}
\renewcommand{\ChapLocation}{Chapters/Chapter_1/}
\begin{absbox}
This is a thesis layout whose example content is its user guide. Start here to build the document, find the files and replace the examples without reorganising your chapters.
\end{absbox}

\section{A working document and a guide}
Read the typeset examples beside their source. A bordered code listing shows
what to type; the figure, equation or paragraph nearby shows how it is used.
Listings are literal text, so a declaration shown in a listing does not define
a command. The actual declarations live in the files named in the explanation.

The design uses the \texttt{memoir} class: coloured chapter openings, chapter
markers on the page edges, running headings and circular page numbers. These
are design choices. The university title pages have a separate role and are
explained in \autoref{chap:02}. This is a community template derived from a thesis,
not an official university-issued class.

\section{Build from the root folder}
Use a recent full TeX Live or MiKTeX installation with \texttt{memoir} and the
packages listed in \texttt{Configuration.tex}. TeX Live 2025 or later is the
intended starting point. Select \texttt{main.tex} as the main document and
pdfLaTeX as the engine. The bibliography backend is BibTeX, selected by
\texttt{backend=bibtex} in the configuration.
\begin{lstlisting}[language={}]
pdflatex main.tex
bibtex main
pdflatex main.tex
pdflatex main.tex
\end{lstlisting}
Run pdfLaTeX again if it still requests a rerun. First-use pages, acronyms and
the table of contents need information from previous runs. An empty symbol list
on the first pass is expected. An alternative is \texttt{latexmk -pdf main.tex},
which manages the required passes. In an online editor, upload the whole folder,
select \texttt{main.tex}, select pdfLaTeX and use its complete rebuild option
after changing bibliography settings.

Build \texttt{propositions.tex} separately with pdfLaTeX. It reads the same
personal details but produces a smaller, independent insert. It is not included
in \texttt{main.tex}. Compilation commands are instructions for the reader;
the distributed source review does not substitute for a build on your system.

\section{Where each piece belongs}
\begin{table}[htbp]
\centering\small
\caption{The existing folder structure and its purpose.}
\label{tab:structure}
\begin{tabularx}{\linewidth}{@{}p{0.43\linewidth}X@{}}
\toprule
Location & Purpose\\\midrule
\texttt{main.tex} & Reading order and document-level switches.\\
\texttt{ThesisDetails.tex} & Shared author, title, date, identifiers and cover paths.\\
\texttt{Configuration.tex} & Packages, page geometry, colours and styles.\\
\texttt{DeclareSymbols.tex} & Global symbol registry and your declarations.\\
\texttt{MathOperators.tex} & Mathematical operators and delimiter helpers.\\
\texttt{Chapters/Chapter\_0} & Covers, title pages, acknowledgements and summaries.\\
\texttt{Chapters/Chapter\_1} to \texttt{Chapter\_8} & One main file per chapter.\\
\texttt{fig/} within each chapter & Figures owned by that chapter.\\
\texttt{Chapters/Chapter\_Appendix} & Acronyms, symbols, notation and back cover.\\
\texttt{References/} & Separate reference and publication databases.\\
\texttt{tests/} & Source checks and an optional typesetting regression document.\\
\bottomrule
\end{tabularx}
\end{table}

\section{Replace content in small steps}
Start with \texttt{ThesisDetails.tex}; then replace the English and Dutch
summaries, acknowledgements and chapter text. Keep figures in their chapter's
\texttt{fig} folder. Empty figure directories contain a \texttt{.gitkeep} file
so Git preserves them. Chapter 8 has no figure directory in the original tree;
add one there only if you later need figures in that chapter.

Retain labels while you are drafting, or update every reference when renaming
them. A consistent pattern such as \texttt{fig:ch3:earth}, \texttt{eq:linear}
and \texttt{sec:acronyms} makes searches useful. Prefer \verb|\autoref{...}|
for named cross-references and \verb|\eqref{...}| for parenthesised equation
numbers. Do not type the resulting numbers by hand.

\section{Source files and generated files}
Share source files, bibliography databases, figure assets and their credits.
Files such as \texttt{.aux}, \texttt{.log}, \texttt{.toc} and \texttt{.bbl} are
build products. The included \texttt{.gitignore} excludes these products while
allowing PDF figures in the chapter directories. It does not ignore every PDF.
The old local copies of standard \texttt{.sty} packages are not bundled; the
template uses the maintained copies in the TeX distribution.
